\documentclass[12pt,a4paper]{report}

% ======================
% PACOTES BÁSICOS
% ======================
\usepackage[utf8]{inputenc}
\usepackage[T1]{fontenc}
\usepackage[brazil]{babel}
\usepackage{geometry}
\usepackage{setspace}
\usepackage{indentfirst}
\usepackage{graphicx}

\usepackage[alf]{abntex2cite} % estilo autor-data (tipo CiteOnLine)
\usepackage{acronym}

\geometry{margin=3cm}
\onehalfspacing

% ======================
% DADOS DO TRABALHO
% ======================
\title{Título do Trabalho}
\author{Seu Nome}
\date{Ano}

% ======================
\begin{document}

% ======================
% CAPA
% ======================
\begin{titlepage}
    \centering
    {\large Universidade Estadual de Maringá\par}
    \vspace{2cm}
    
    {\Large\textbf{Vitor Azevedo Padovani} \par}
    \vspace{2cm}
    
    {\LARGE\textbf{Análise de Logs de Rede para Detecção de Anomalias Utilizando Aprendizado de Máquina: Uma Extensão com Engenharia de Atributos e Avaliação Comparativa Aprofundada} \par}
    \vspace{2cm}
    
    {\large Trabalho de Conclusão de Curso \par}
    \vspace{1cm}
    
    {\large Orientadora: Luciana Andréia Fondazzi Martimiano \par}
    \vfill
    
    {\large Maringá \par}
    {\large 2026 \par}
\end{titlepage}

% ======================
% RESUMO
% ======================
\chapter*{Resumo}
\textit{Logs} são registros gerados continuamente durante a operação de equipamentos e sistemas,
incluindo dispositivos de rede, servidores web, sistemas operacionais e aplicações. Sua importância
está em auditar ações de usuários, identificar comportamentos recorrentes e detectar conexões
incomuns, como àquelas iniciadas por ferramentas maliciosas, \textit{scanners} automatizados ou usuários
mal-intencionados, as quais são caracterizadas como comportamentos anômalos. Embora
indispensáveis para a segurança da informação, \textit{logs} são de difícil análise manual: possuem formato
de texto não estruturado, variam conforme o fabricante e podem gerar centenas de milhares de
registros por hora. Essa situação é encontrada na administração do NPD (Núcleo de Processamento
de Dados) da UEM (Universidade Estadual de Maringá), onde os \textit{logs} de acesso do servidor web
Nginx registram o tráfego de dezenas de sites institucionais. Este trabalho propõe uma extensão do
trabalho de Maziero (2024), que aplicou os algoritmos Isolation Forest, Local Outlier Factor e
Robust Covariance sobre os \textit{logs} do NPD, utilizando seis atributos com codificação manual e
avaliando os modelos com dez anomalias inseridas artificialmente. As principais limitações
identificadas naquela abordagem são: a perda de informação semântica pelo \textit{encoding} simplificado,
a ausência de \textit{features} comportamentais temporais, e a avaliação com dados sintéticos em vez de
padrões reais do ambiente. Este trabalho pretende superar essas limitações por meio de três
contribuições principais: primeiro, uma engenharia de atributos estruturada em quatro famílias:
atributos da URI (Uniform Resource Identifier), atributos temporais com codificação cíclica,
atributos comportamentais por IP (Internet Protocol) em janela deslizante, atributos do \textit{User Agent};
segundo, a construção de um \textit{ground truth} baseado em evidências reais do ambiente, em substituição
às anomalias sintéticas da abordagem anterior; terceiro, uma avaliação comparativa entre Isolation
Forest e One-Class SVM utilizando métricas adequadas para \textit{datasets} desbalanceados.

\textbf{Palavras-chave:} Gerenciamento de Redes. Detecção de anomalias. Aprendizado de Máquina. \textit{Logs}.

% ======================
% SUMÁRIO
% ======================
\tableofcontents
\newpage

% ======================
% SIGLAS E ABREVIAÇÕES
% ======================
\chapter*{Lista de Siglas e Abreviaturas}
\begin{acronym}
	\acro{IA}{Inteligência Artificial}
	\acro{ML}{Machine Learning}
	\acro{NPD}{Núcleo de Processamento de Dados}
	\acro{UEM}{Universidade Estadual de Maringá}
	\acro{URI}{Uniform Resource Identifier}
	\acro{IP}{Internet Protocol}
	\acro{NIST}{National Institute of Standards and Technology}
	\acro{CID}{Confidencialidade, Integridade e Disponibilidade}
	\acro{LGPD}{Lei Geral de Proteção de Dados}
	\acro{MITM}{Man-in-the-Middle}
	\acro{ARP}{Address Resolution Protocol}
	\acro{MAC}{Media Access Control}
	\acro{IDS}{Intrusion Detection System}
	\acro{AAA}{Authentication, Authorization and Accountability}
	\acro{VPN}{Virtual Private Network}
	\acro{MFA}{Multi-Factor Authentication}
	\acro{SIEM}{Security Information and Event Management}
	\acro{SVM}{Support Vector Machine}
\end{acronym}

% ======================
% INTRODUÇÃO
% ======================
\chapter*{Introdução}

Com a crescente digitalização de serviços e o uso de sistemas interconectados, a segurança computacional tornou-se um desafio central para organizações. Aplicações web, em especial, apresentam alta exposição a ameaças, sendo frequentemente alvo de ataques automatizados e exploração de vulnerabilidades. Relatórios recentes indicam aumento contínuo na frequência e complexidade desses ataques, impulsionado, entre outros fatores, pelo uso de inteligência artificial por agentes maliciosos \cite{forbes2024}. Esse cenário reforça a necessidade de mecanismos eficazes de monitoramento e detecção de comportamentos anômalos.

Nesse contexto, logs de rede são fundamentais, pois registram eventos detalhados sobre sistemas e usuários, sendo a principal fonte para auditoria e detecção de incidentes. Logs de servidores web, como o Nginx, incluem informações como requisições HTTP, endereços IP e códigos de resposta, permitindo reconstruir interações com o sistema. Estudos mostram que sua análise é essencial para identificar atividades maliciosas \cite{hussein2024, komadina2024}.

Entretanto, a análise de logs apresenta desafios relevantes. Grandes volumes de dados são gerados continuamente, muitas vezes em formatos não estruturados, dificultando a interpretação manual. Além disso, anomalias tendem a ser raras e sutis, ocultas em meio a comportamentos normais \cite{hussein2024}. Essa combinação torna inviável a análise manual, exigindo abordagens automatizadas.

Diante desse cenário, técnicas de aprendizado de máquina têm se consolidado como uma abordagem promissora para a detecção de anomalias em logs. A detecção de anomalias consiste em identificar padrões que se desviam significativamente do comportamento esperado, podendo indicar falhas, erros ou atividades maliciosas. Métodos não supervisionados, como Isolation Forest e One-Class SVM, são particularmente relevantes nesse domínio, pois não dependem de dados previamente rotulados, característica importante em cenários reais onde a rotulação de ataques é limitada ou inexistente. Esses métodos são capazes de identificar tanto ataques conhecidos quanto comportamentos inéditos, contribuindo para a detecção de ameaças emergentes \cite{komadina2024}.

Apesar dos avanços na área, a eficácia desses modelos está diretamente relacionada à qualidade da representação dos dados utilizada como entrada. Trabalhos recentes apontam que muitos métodos utilizam representações simplificadas dos logs, desconsiderando informações semânticas e contextuais relevantes, como padrões de acesso ao longo do tempo ou características detalhadas das requisições. Essa limitação pode comprometer a capacidade dos modelos de capturar anomalias mais complexas, especialmente aquelas relacionadas ao comportamento de usuários ou à exploração de recursos específicos do sistema \cite{chen2024semantic}.

Nesse contexto, o trabalho de Maziero (2024) representa uma contribuição inicial ao aplicar algoritmos de aprendizado de máquina não supervisionado na detecção de anomalias em logs reais do servidor Nginx da Universidade Estadual de Maringá (UEM). Contudo, sua abordagem baseia-se em um conjunto restrito de atributos, o que limita o aproveitamento das informações disponíveis nos logs.

Dessa forma, o presente trabalho tem como objetivo geral desenvolver um pipeline de detecção de anomalias em logs de rede com foco no enriquecimento da representação dos dados por meio de engenharia de atributos. Especificamente, busca-se incorporar informações semânticas da URI, padrões temporais de comportamento por endereço IP e características do \textit{User Agent}, avaliando o impacto dessas informações no desempenho dos algoritmos Isolation Forest e One-Class SVM.

A motivação deste trabalho está na necessidade de aprimorar a gestão de logs e fortalecer a segurança da infraestrutura computacional da UEM. Ao propor uma abordagem mais rica em termos de representação de dados e mais rigorosa em termos de avaliação, espera-se contribuir para a detecção mais precisa de anomalias. Além disso, a solução proposta possui potencial de replicação em outros ambientes institucionais, contribuindo para o avanço da aplicação de técnicas de aprendizado de máquina na área de segurança da informação.

% ======================
% REFERENCIAL TEÓRICO
% ======================
\chapter{Fundamentação Teórica}

Este capítulo apresenta os conceitos que fundamentam o desenvolvimento deste trabalho. Os temas cobertos foram selecionados para situar o problema de detecção de anomalias em logs de rede dentro dos campos da segurança da informação, do gerenciamento de redes e da inteligência artificial.

% ─────────────────────────────────────────────────────────────────────────────
\section{Segurança da Informação}
% ─────────────────────────────────────────────────────────────────────────────

Toda informação que transita por uma rede, seja por meios físicos ou
digitais, precisa de mecanismos que assegurem que ela chegue ao destino
sem alterações, seja legível apenas para quem tem direito a isso, e
permaneça acessível quando necessário. Esse conjunto de preocupações
constitui o domínio da segurança da informação.

Do ponto de vista técnico, segurança da informação envolve a definição de
métodos e parâmetros capazes de reduzir ao mínimo
os riscos a que ativos computacionais estão
expostos \cite{stallings2014}. Além de proteger dados em trânsito, a
disciplina se preocupa com a continuidade dos negócios, a minimização de
riscos e a proteção dos investimentos organizacionais, o que exige como
primeiro passo uma análise cuidadosa do que deve ser protegido e de quais
ameaças precisam ser consideradas \cite{hintzbergen2018}.

Para padronizar esses requisitos entre diferentes sistemas e
organizações, o \textit{National Institute of Standards and Technology}
(NIST) publicou diretrizes para categorização de informações e sistemas
de informação~\cite{nist2004}, a partir das quais a segurança da
informação passou a ser descrita por três pilares fundamentais:
Confidencialidade, Integridade e Disponibilidade (CID), dos quais se
derivam os demais conceitos da área. Complementarmente,
\citeonline{kurose2021} identificam quatro propriedades essenciais de uma
comunicação segura: confidencialidade, integridade de mensagem,
autenticação de ponto final e segurança operacional.

\subsection{Confidencialidade}

Confidencialidade é a propriedade que garante que uma mensagem só pode
ser lida por aqueles a quem ela se destina. Do ponto de vista do NIST,
trata-se de preservar restrições de acesso e divulgação, protegendo a
privacidade pessoal e informações proprietárias, sendo a perda de
confidencialidade caracterizada pela revelação não autorizada de
informações~\cite{nist2004}.

O conceito se desdobra em dois aspectos complementares: a
\textit{confidencialidade de dados}, que restringe a leitura de
informações privadas aos indivíduos autorizados; e a
\textit{privacidade}, que devolve ao titular dos dados o controle sobre
quem pode coletá-los, armazená-los e acessá-los~\cite{stallings2014}.
Essa proteção deve se estender por todo o ciclo de vida da informação,
enquanto ela estiver armazenada em sistemas, em trânsito pela rede e ao
chegar ao destino~\cite{hintzbergen2018}.

Um exemplo concreto do impacto da confidencialidade no cenário nacional é
a LGPD (Lei Geral de Proteção de Dados), que responsabiliza as
organizações pelos dados pessoais de brasileiros que processam. Quando
uma organização não consegue impedir o acesso não autorizado a esses
dados, independentemente de culpa, a confidencialidade foi violada e
consequências legais se tornam aplicáveis~\cite{serpro2024}.

A criptografia é o principal mecanismo técnico para garantir essa
propriedade: ao cifrar uma mensagem com um método robusto, garante-se que
um eventual interceptador não consiga extrair seu
conteúdo~\cite{kurose2021}.

\subsection{Integridade}

Se a confidencialidade protege quem pode ler uma informação, a
integridade protege o conteúdo dessa informação contra modificações não
autorizadas. A garantia é a de que a mensagem recebida é exatamente
aquela que foi enviada, sem alterações ao longo do caminho.

O NIST define integridade como a salvaguarda da informação contra
modificação ou destruição imprópria, incluindo a garantia de
autenticidade e não repúdio. A perda dessa propriedade se caracteriza
pela destruição ou alteração não autorizada de dados~\cite{nist2004}.
Segundo \citeonline{stallings2014}, a integridade abrange dois níveis: a
\textit{integridade de dados}, que determina que informações e programas
só sejam alterados por processos e pessoas autorizadas; e a
\textit{integridade de sistema}, que exige que um programa ou sistema
execute suas funções previstas de forma inalterada ao longo de toda sua
vida útil.

Um exemplo clássico de ataque a essa propriedade é o
\textit{Man-in-the-Middle} (MITM) por \textit{ARP Cache Poisoning}. O
protocolo ARP traduz endereços IP, da camada de rede, para endereços MAC,
na camada de enlace. No ataque, as respostas ARP são falsificadas para
redirecionar o tráfego por uma máquina controlada pelo atacante, que
passa a interceptar a comunicação. O resultado é a violação simultânea da
integridade, pela adulteração das respostas, e da confidencialidade,
pelo acesso indevido às informações trocadas \cite{fortinet2024}.
Tecnicamente, a integridade pode ser verificada por meio de somas de
verificação e funções de \textit{hash} criptográfico aplicadas nas
camadas de enlace e transporte~\cite{kurose2021}.

\subsection{Disponibilidade}

Disponibilidade é a capacidade de acessar um sistema e seus serviços
sempre que necessário, dentro do período de operação previsto, mesmo
diante de falhas ou tentativas de interrupção. Mais do que simplesmente
estar \textit{online}, um sistema disponível deve manter sua
operacionalidade mesmo sob condições adversas.

Sistemas críticos ilustram bem as consequências da indisponibilidade:
infraestruturas hospitalares, controle de tráfego aéreo, redes elétricas
e usinas nucleares não toleram interrupções. Estimativas indicam que até
2025 serão produzidos 79,4 zettabytes de dados, com cerca de 30\% deles
exigindo processamento por sistemas de tempo real~\cite{intel2024},
evidenciando a crescente demanda por ambientes resilientes.

O NIST define disponibilidade como a garantia do acesso oportuno e
confiável à informação e aos sistemas, sendo sua perda a interrupção
desse acesso~\cite{nist2004}. \citeonline{hintzbergen2018} decompõe essa
propriedade em três características: \textit{oportunidade} (informações
acessíveis quando necessário), \textit{continuidade} (manutenção do
serviço mesmo em caso de falha) e \textit{robustez} (capacidade de
atender todos os usuários autorizados simultaneamente).
Na prática, a disponibilidade é suportada por \textit{firewalls} e
sistemas de detecção de intrusão (IDS), que monitoram anomalias no
tráfego e intervêm antes que ataques de negação de serviço comprometam
a operação~\cite{kurose2021}.

\subsection{Autorização, Autenticação e Não Repúdio}

A tríade CID define \textit{o que} deve ser protegido, mas não aborda
\textit{quem} acessa os recursos nem \textit{o que} cada pessoa pode
fazer. Para cobrir essa lacuna, o modelo AAA (\textit{Authentication,
	Authorization and Accountability}) estabelece três propriedades
complementares~\cite{convery2007}.

\subsubsection{Autenticação}

Autenticação responde à pergunta \textit{"Quem é você?"} \cite{
	convery2007}. É o mecanismo que verifica se um usuário é de fato quem
alega ser antes de conceder qualquer acesso. Em um cenário de acesso
remoto via VPN, por exemplo, a autenticação por senha ou por múltiplos
fatores (MFA) impede que credenciais roubadas sejam suficientes para
comprometer o sistema.

\subsubsection{Autorização}

Autorização responde à pergunta \textit{"O que você pode fazer?"} \cite{
	convery2007}. Mesmo após autenticado, um usuário deve operar apenas dentro
dos limites do seu perfil. Um estagiário com acesso integral à
infraestrutura de rede representa um risco latente, não pela intenção,
mas pela possibilidade de consequências graves a partir de ações
cotidianas fora do escopo de sua função. O princípio do
\textit{Least Privilege}, conceder apenas os privilégios estritamente
necessários, é a forma mais direta de mitigar esse risco.

\subsubsection{Não Repúdio}

Não repúdio responde à pergunta \textit{"O que você fez?"}~\cite{
	convery2007}. Sem registros rastreáveis das ações realizadas em um
ambiente, incidentes podem ocorrer sem que seja possível identificar
origem, causa ou responsável. A análise de logs de sistemas e
dispositivos de rede é o instrumento principal para garantir essa
propriedade: registros de usuários autenticados, endereços IP, horários
de acesso e operações realizadas formam uma trilha de auditoria que
impede a evasão de responsabilidade. Quando integrada a
\textit{firewalls}, IDS e SIEMs (\textit{Security Information and Event
	Management}), essa análise torna-se proativa, permitindo identificar
comportamentos suspeitos antes que causem danos~\cite{convery2007}.

% ─────────────────────────────────────────────────────────────────────────────
\section{Gerenciamento de Redes}
% ─────────────────────────────────────────────────────────────────────────────

...

% ======================
% REFERÊNCIAS
% ======================
\bibliography{references} 

\end{document}